\documentclass[11pt,a4paper]{article}

\usepackage[utf8]{inputenc}
\usepackage[T1]{fontenc}
\usepackage[french]{babel}
\usepackage[a4paper,margin=2.1cm]{geometry}
\usepackage{amsmath,amssymb,amsthm,mathtools}
\usepackage{lmodern}
\usepackage{enumitem}
\usepackage{hyperref}
\usepackage{xcolor}
\usepackage{fancyhdr}
\usepackage{lastpage}
\usepackage{tikz}
\usepackage{stmaryrd}
\usepackage{array}
\usepackage{tabularx}

\hypersetup{
    colorlinks=true,
    linkcolor=black,
    urlcolor=blue
}

\setlength{\parindent}{0pt}
\setlength{\parskip}{0.6em}

\pagestyle{fancy}
\fancyhf{}
\fancyhead[L]{Devoir surveillé --- Nombres complexes}
\fancyhead[R]{Prépa scientifique}
\fancyfoot[C]{\thepage/\pageref{LastPage}}

\newtheorem*{remarque}{Remarque}
\newtheorem*{consignes}{Consignes générales}

\newcommand{\C}{\mathbb{C}}
\newcommand{\R}{\mathbb{R}}
\newcommand{\U}{\mathbb{U}}
\newcommand{\ii}{\mathrm{i}}
\newcommand{\e}{\mathrm{e}}
\newcommand{\dd}{\,\mathrm{d}}

\newcommand{\bareme}[1]{\hfill{\small\textbf{[#1 pts]}}}
\newcommand{\ds}{\displaystyle}

\begin{document}

\begin{center}
    {\LARGE \textbf{Devoir surveillé de mathématiques}}\\[0.4em]
    {\Large \textbf{Nombres complexes}}\\[0.8em]
    \rule{0.9\linewidth}{0.6pt}\\[0.5em]
    \textbf{Niveau :} début de prépa scientifique \qquad
    \textbf{Durée conseillée :} 4 heures \qquad
    \textbf{Calculatrice interdite}
\end{center}

\begin{consignes}
La qualité de la rédaction, la précision des justifications et la clarté des raisonnements seront prises en compte de façon importante dans la notation.

Toute réponse non justifiée pourra être très faiblement valorisée, même si le résultat annoncé est exact.

Aucune initiative pertinente ne sera pénalisée, même si elle n’aboutit pas complètement.

Le barème proposé dans ce sujet est \textbf{indicatif}. Il valorise davantage les questions demandant une idée de fond, une rédaction plus longue, une méthode non immédiate ou une vraie maîtrise globale du cours.
\end{consignes}

\vspace{0.6em}

\begin{center}
\renewcommand{\arraystretch}{1.25}
\begin{tabular}{|>{\centering\arraybackslash}m{4.2cm}|>{\centering\arraybackslash}m{3.2cm}|>{\centering\arraybackslash}m{6.3cm}|}
\hline
\textbf{Partie} & \textbf{Points indicatifs} & \textbf{Idée dominante} \\
\hline
Exercice 1 & 46 pts & Racines de l'unité, polygones réguliers, barycentre, structure algébrique et géométrique \\
\hline
Exercice 2 & 18 pts & Transformation homographique, cercle unité, interprétation géométrique \\
\hline
Exercice 3 & 20 pts & Géométrie complexe, construction d'un carré, rotation complexe \\
\hline
Exercice 4 & 8 pts & c'est qui lui? \\
\hline
\textbf{Total indicatif} & \textbf{94 pts} & \\
\hline
\end{tabular}
\end{center}

\bigskip

\section*{Rappels utiles}

Les rappels suivants ne dispensent pas d'une rédaction rigoureuse. Ils visent seulement à remettre en mémoire
quelques outils qui pourront être utilisés dans le sujet.

\subsection*{1. Somme des termes d'une suite géométrique}

Si \(q\neq 1\), on a
\[
1+q+q^2+\cdots+q^{n-1}=\sum_{k=0}^{n-1}q^k=\frac{1-q^n}{1-q}.
\]
En particulier, si \(q^n=1\) et \(q\neq 1\), alors
\[
\sum_{k=0}^{n-1}q^k=0.
\]

\subsection*{2. Écriture d'un polynôme dans \(\C[X]\)}

Dire que \(P\in \C[X]\) signifie que \(P\) est un polynôme à coefficients complexes. On peut donc l'écrire sous la forme
\[
P(X)=a_0+a_1X+\cdots+a_mX^m,
\qquad a_0,\dots,a_m\in\C.
\]
Si \(P\in \C_{n-1}[X]\), cela signifie que \(\deg(P)\le n-1\), donc
\[
P(X)=a_0+a_1X+\cdots+a_{n-1}X^{n-1}.
\]
Cette écriture permet souvent de traiter \(P\) terme à terme.

\subsection*{3. Quand peut-on intervertir une somme double ?}

Pour des sommes \emph{finies}, on peut toujours changer l'ordre de sommation :
\[
\sum_{\omega\in E}\sum_{k=0}^{n-1} a_{\omega,k}
=
\sum_{k=0}^{n-1}\sum_{\omega\in E} a_{\omega,k}.
\]
Cela vient du fait qu'une somme finie n'est qu'une addition d'un nombre fini de termes, et que l'addition est associative et commutative.

\subsection*{4. Notation produit \(\prod\)}

La notation
\[
\prod_{k=0}^{n-1} u_k
\]
désigne le produit
\[
u_0u_1\cdots u_{n-1}.
\]
De même,
\[
\prod_{\omega\in \U_n}(X-z_\omega)
\]
désigne le produit de tous les facteurs \(X-z_\omega\) lorsque \(\omega\) parcourt l'ensemble \(\U_n\).

\subsection*{5. Polygone régulier}

Un polygone est dit \emph{régulier} lorsque tous ses côtés ont la même longueur et que ses sommets sont répartis de manière uniforme sur un même cercle. En termes complexes, les racines \(n\)-ièmes de l'unité
\[
1,\e^{2\ii\pi/n},\e^{4\ii\pi/n},\dots,\e^{2(n-1)\ii\pi/n}
\]
forment les sommets d'un polygone régulier à \(n\) côtés centré en \(0\).

\subsection*{6. Barycentre}

Si \(z_1,\dots,z_n\) sont les affixes de points \(M_1,\dots,M_n\), le barycentre de ces points affectés de coefficients égaux est le point \(G\) d'affixe
\[
z_G=\frac{z_1+\cdots+z_n}{n}.
\]
Autrement dit, la moyenne des affixes donne l'affixe du ``centre moyen'' de la configuration.
Pour les sommets d'un polygone régulier, ce barycentre est le centre du polygone.

\subsection*{7. Rappel sur la cotangente}

Pour tout réel \(\theta\) tel que \(\sin\theta\neq 0\), on définit la cotangente de \(\theta\) par
\[
\cot(\theta)=\frac{\cos\theta}{\sin\theta}.
\]
En particulier,
\[
\cot\!\left(\frac{\theta}{2}\right)=\frac{\cos(\theta/2)}{\sin(\theta/2)}
\]
dès que \(\sin(\theta/2)\neq 0\), c'est-à-dire dès que \(\theta\notin 2\pi\mathbb{Z}\).

Cette fonction apparaît naturellement lorsque l'on simplifie des quotients du type
\[
\frac{1+\e^{\ii\theta}}{1-\e^{\ii\theta}}.
\]
\subsection*{8. Demi-plan supérieur}

On appelle \emph{demi-plan supérieur} l'ensemble
\[
\{z\in\C\mid \Im(z)>0\}.
\]
Si \(z=x+\ii y\), cela signifie simplement que \(y>0\).
Géométriquement, c'est l'ensemble des points du plan complexe situés strictement au-dessus de l'axe réel.

\subsection*{9. Que signifie ``stable par'' ?}

Soit \(E\) un ensemble muni d'une opération.

Dire que \(E\) est \emph{stable par} cette opération signifie que lorsque l'on applique cette opération à des éléments de \(E\), on obtient encore un élément de \(E\).

Par exemple :
\begin{itemize}
    \item dire que \(E\) est stable par addition signifie que
    \[
    \forall x,y\in E,\quad x+y\in E;
    \]
    \item dire que \(E\) est stable par produit signifie que
    \[
    \forall x,y\in E,\quad xy\in E;
    \]
    \item dire qu'un ensemble de nombres non nuls est stable par passage à l'inverse signifie que
    \[
    \forall x\in E,\quad x\neq 0 \Longrightarrow x^{-1}\in E.
    \]
\end{itemize}

Ainsi, montrer que \(\U_n\) est stable par produit et par passage à l'inverse revient à montrer que :
\[
\forall \omega,\omega'\in\U_n,\quad \omega\omega'\in\U_n
\qquad\text{et}\qquad
\forall \omega\in\U_n,\quad \omega^{-1}\in\U_n.
\]
\bigskip

\section*{Exercice 1 --- Racines de l'unité, barycentre complexe et polygones réguliers \hfill \textbf{46 pts}}

On note
\[
\U_n=\left\{\omega\in\C\mid \omega^n=1\right\},
\]
où \(n\ge 2\) est un entier fixé.

\subsection*{Partie A --- Premières propriétés des racines \(n\)-ièmes de l'unité \hfill \textbf{11 pts}}

\begin{enumerate}[label=\textbf{A.\arabic*.}]
    \item Déterminer explicitement les éléments de \(\U_n\).
    \bareme{2}

    \item Montrer que
    \[
    \U_n=\left\{\e^{2\ii k\pi/n}\mid k\in\llbracket 0,n-1\rrbracket\right\}.
    \]
    \bareme{1.5}

    \item Montrer que \(\U_n\) est stable par produit et par passage à l'inverse.
    \bareme{1.5}

    \item Calculer
    \[
    \sum_{\omega\in\U_n}\omega.
    \]
    \bareme{2}

    \item Plus généralement, pour \(m\in\mathbb{Z}\), calculer
    \[
    S_m=\sum_{\omega\in\U_n}\omega^m.
    \]
    \bareme{2.5}

    \item En déduire que, pour tout polynôme \(P\in\C_{n-1}[X]\),
    \[
    \frac1n\sum_{\omega\in\U_n}P(\omega)=P(0).
    \]
    \bareme{1.5}
\end{enumerate}

\subsection*{Partie B --- Du cercle unité au polygone régulier \hfill \textbf{12 pts}}

On considère un complexe \(a\in\C\) et un réel \(\rho>0\). Pour tout \(\omega\in\U_n\), on pose
\[
z_\omega=a+\rho\,\omega.
\]

\begin{enumerate}[label=\textbf{B.\arabic*.}]
    \item Montrer que les points d'affixes \(z_\omega\) appartiennent tous à un même cercle, dont on précisera le centre et le rayon.
    \bareme{1.5}

    \item Montrer que, pour \(\omega\neq \omega'\),
    \[
    |z_\omega-z_{\omega'}|=\rho\,|\omega-\omega'|.
    \]
    \bareme{1.5}

    \item En déduire que les points d'affixes \(z_\omega\) sont les sommets d'un polygone régulier à \(n\) côtés.
    \bareme{2.5}

    \item Montrer que
    \[
    \frac1n\sum_{\omega\in\U_n} z_\omega=a.
    \]
    \bareme{2}

    \item Interpréter géométriquement le résultat précédent.
    \bareme{1}

    \item Montrer que
    \[
    \prod_{\omega\in\U_n}(X-z_\omega)=(X-a)^n-\rho^n.
    \]
    \bareme{2.5}

    \item Retrouver à partir de cette identité que les \(z_\omega\) sont les sommets d'un polygone régulier.
    \bareme{1}
\end{enumerate}

\subsection*{Partie C --- Réciproque partielle \hfill \textbf{11 pts}}

On se donne \(n\) complexes distincts \(z_0,\dots,z_{n-1}\). On suppose qu'il existe \(a\in\C\) et \(\lambda\in\C^\ast\) tels que
\[
z_k-a=\lambda\,\e^{2\ii k\pi/n}
\qquad (k\in\llbracket 0,n-1\rrbracket).
\]

\begin{enumerate}[label=\textbf{C.\arabic*.}]
    \item Montrer que tous les points d'affixes \(z_k\) sont situés sur un même cercle.
    \bareme{1.5}

    \item Montrer que
    \[
    \sum_{k=0}^{n-1} z_k=na.
    \]
    \bareme{1.5}

    \item Montrer que, pour tout entier \(m\),
    \[
    \sum_{k=0}^{n-1}(z_k-a)^m=
    \begin{cases}
    0 & \text{si } n\nmid m,\\[0.4em]
    n\lambda^m & \text{si } n\mid m.
    \end{cases}
    \]
    \bareme{3.5}

    \item Montrer que, pour tout polynôme \(P\in\C_{n-1}[X]\),
    \[
    \frac1n\sum_{k=0}^{n-1}P(z_k)=P(a).
    \]
    \bareme{3}

    \item Donner une interprétation géométrique de cette propriété.
    \bareme{1.5}
\end{enumerate}

\subsection*{Partie D --- Cas particuliers classiques \hfill \textbf{12 pts}}

\subsubsection*{1. Triangle équilatéral \hfill \textbf{6.5 pts}}

Soient \(z_1,z_2,z_3\) trois complexes distincts de même module.

\begin{enumerate}[label=\textbf{D1.\arabic*.}]
    \item Montrer que si
    \[
    z_1+z_2+z_3=0,
    \]
    alors \(z_1,z_2,z_3\) sont les sommets d'un triangle équilatéral.
    \bareme{2.5}

    \item Montrer la réciproque suivante : si \(z_1,z_2,z_3\) sont les sommets d'un triangle équilatéral centré en \(0\), alors
    \[
    z_1+z_2+z_3=0.
    \]
    \bareme{2}

    \item Montrer qu'un triangle équilatéral quelconque admet, quitte à renuméroter ses sommets, une écriture de la forme
    \[
    z_1=a+r\e^{\ii\theta},\qquad
    z_2=a+r\e^{\ii(\theta+2\pi/3)},\qquad
    z_3=a+r\e^{\ii(\theta+4\pi/3)},
    \]
    avec \(a\in\C\), \(r>0\) et \(\theta\in\R\).
    \bareme{2}
\end{enumerate}

\subsubsection*{2. Carré \hfill \textbf{5.5 pts}}

Soient \(z_1,z_2,z_3,z_4\) quatre complexes distincts de même module.

\begin{enumerate}[label=\textbf{D2.\arabic*.}]
    \item Montrer que si
    \[
    z_1+z_2+z_3+z_4=0,
    \]
    alors, après renumérotation éventuelle, ces points sont les sommets d'un rectangle centré en \(0\).
    \bareme{2.5}

    \item Donner une condition supplémentaire simple pour que ce rectangle soit un carré.
    \bareme{1.5}

    \item Traduire cette condition sous forme complexe.
    \bareme{1.5}
\end{enumerate}

\bigskip

\section*{Exercice 2 --- Cercle unité et transformation homographique \hfill \textbf{18 pts}}

Soit \(u\in\C\) tel que \(|u|=1\) et \(u\neq 1\). On pose
\[
w=\frac{1+u}{1-u}.
\]

\begin{enumerate}[label=\textbf{\arabic*.}]
    \item Montrer que
    \[
    \overline w=-w.
    \]
    \bareme{2.5}

    \item En déduire que \(w\in \ii\R\).
    \bareme{1}

    \item Réciproquement, soit \(w\in \ii\R\). Montrer que
    \[
    u=\frac{w-1}{w+1}
    \]
    est bien défini et vérifie \(|u|=1\).
    \bareme{3}

    \item Montrer que l'application
    \[
    u\longmapsto \frac{1+u}{1-u}
    \]
    réalise une bijection du cercle unité privé du point \(1\) vers l'axe imaginaire pur.
    \bareme{3}

    \item Soit \(\theta\in\R\setminus 2\pi\mathbb{Z}\). En prenant \(u=\e^{\ii\theta}\), montrer que
    \[
    \frac{1+\e^{\ii\theta}}{1-\e^{\ii\theta}}
    =\ii\,\cot\!\left(\frac{\theta}{2}\right).
    \]
    \bareme{3}

    \item En déduire une expression de
    \[
    \cot\!\left(\frac{\theta}{2}\right)
    \]
    en fonction de \(\e^{\ii\theta}\).
    \bareme{1}

    \item Résoudre dans \([0,2\pi[\) l'équation
    \[
    \frac{1+\e^{\ii\theta}}{1-\e^{\ii\theta}}=\ii\sqrt3.
    \]
    \bareme{1.5}

    \item Soit \(z\in\C\) avec \(z\neq 1\). Montrer que les deux propriétés suivantes sont équivalentes :
    \[
    |z|=1
    \qquad\Longleftrightarrow\qquad
    \frac{1+z}{1-z}\in \ii\R.
    \]
    \bareme{3}

    \item Interpréter géométriquement ce résultat : expliquer pourquoi cette transformation envoie un cercle sur une droite.
    \bareme{2}
\end{enumerate}

\bigskip

\section*{Exercice 3 --- Un carré qui tourne \hfill \textbf{20 pts}}

Dans le plan complexe rapporté à un repère orthonormé direct, on considère deux points distincts \(A\) et \(B\) d'affixes respectives \(a\) et \(b\).

On cherche à construire les deux carrés ayant \([AB]\) pour côté.

\begin{enumerate}[label=\textbf{\arabic*.}]
    \item Soit \(u=b-a\). Expliquer pourquoi multiplier \(u\) par \(\ii\) revient à effectuer une rotation d'angle \(\pi/2\).
    \bareme{2}

    \item Montrer que si \(C\) est le sommet d'un carré direct \(ABCD\), alors
    \[
    z_C-z_B=\ii\,(b-a).
    \]
    En déduire l'affixe de \(C\).
    \bareme{3}

    \item Montrer que l'autre carré construit sur \([AB]\) correspond à
    \[
    z_C-z_B=-\ii\,(b-a).
    \]
    Donner alors les deux affixes possibles de \(C\).
    \bareme{2}

    \item En déduire les deux affixes possibles de \(D\).
    \bareme{2}

    \item Vérifier directement, dans chacun des deux cas, que les quatre côtés ont même longueur et que deux côtés consécutifs sont perpendiculaires.
    \bareme{3}

    \item On suppose maintenant
    \[
    a=1+\ii,\qquad b=4+2\ii.
    \]
    Calculer explicitement les affixes des deux points \(C\) possibles et des deux points \(D\) correspondants.
    \bareme{2.5}

    \item Déterminer l'affixe du centre \(G\) de chacun de ces carrés.
    \bareme{1.5}

    \item Montrer que les deux centres obtenus sont situés sur la médiatrice de \([AB]\).
    \bareme{2}

    \item Plus généralement, montrer qu'une rotation de centre \(a\) et d'angle \(\theta\) s'écrit
    \[
    z'-a=\e^{\ii\theta}(z-a).
    \]
    \bareme{2}
\end{enumerate}

\bigskip


\section*{Exercice 4 --- Mais c'est qui ce nombre il est imaginaire ?? \hfill \textbf{16 pts}}

On part à l'exploration du plan complexe à la recherche d'un individu très particulier.
Les indices s'accumulent, les pistes se croisent, et pourtant le suspect semble toujours vouloir se cacher.

On cherche donc à déterminer le complexe \(z\in\C\setminus\{-1,1\}\) vérifiant simultanément les propriétés suivantes :

\begin{itemize}
    \item le point d'affixe \(z\) appartient au cercle unité ;
    \item le quotient
    \[
    \frac{1+z}{1-z}
    \]
    est un imaginaire pur ;
    \item le quotient
    \[
    \frac{z-\ii}{z+\ii}
    \]
    est réel ;
    \item enfin, \(z\) appartient au demi-plan supérieur :
    \[
    \Im(z)>0.
    \]
\end{itemize}

Déterminer \(z\).

\bareme{8}

\end{enumerate}

\vfill

\begin{center}
    \rule{0.9\linewidth}{0.6pt}\\[0.4em]
    \textit{Fin du sujet}
\end{center}

\end{document}